\begin{textsample}{NEG Dim 8 – global\_north\_2024\_05 – Score -1.00 – global\_north\_2024\_05/108584383.txt}  \label{ex:f8_neg_274}
Fox News Digital has obtained satellite images showing the now-destroyed temporary pier the U.S. military installed on the coast of Gaza earlier this month.
The \$320 million structure lasted roughly a week before choppy weather battered it apart.
President Biden ' s administration says it is working to repair the pier, but they offered no timeline for completion.
The pier had been used as a \textbf{conduit} for delivering aid to Gaza.
The images show that less than a third of the pier remains intact, and there is no sign of the remnants of the deeper sections of the structure.
The U.S. announced that it was forced to suspend deliveries of aid to the pier on Tuesday, though much of the damage occurred prior to then.
Of the four stabilizing vessels that detached earlier this week, two of the boats floated northward and landed on a beach in Ashdod, Israel, while the two others remain anchored at the beach near the pier.
While the pier has been used to transfer roughly 569 metric tons of aid into Gaza, none of that @ @ @ @ @ @ @ @ @ @, the Pentagon confirmed.
While the pier has been used to transfer roughly 569 metric tons of aid into Gaza, none of that aid had been delivered to Palestinians as of last week, the Pentagon confirmed.
The pier ' s failure comes as Israel conducted a sizable operation in Rafah, with tanks rolling into the heart of the city for the first time since the war began.
Of the four stabilizing vessels that detached earlier this week, two of the boats floated northward and landed on a beach in Ashdod, Israel, while the two others remain anchored at the beach near the pier. ( U.S.
Army via AP ) Witnesses in Rafah told Reuters the Israeli military appeared to be using remote-operated armored vehicles, as there was no immediate sign of personnel in or around them.
The Israeli Defense Forces did not comment on those reports.
Anders Hagstrom is a reporter with Fox News Digital covering national politics and major breaking news events.
Send tips to Anders.Hagstrom@Fox.com, or on Twitter : @ @ @ @ @ @ @ @ @ @ Senior Data reporter Harry Enten revealed numbers on the network Wednesday showing that Americans ' opinions on former President Trump haven't changed during his ongoing hush money trial in New York City.
CNN senior data reporter Harry Enten revealed numbers on the network Wednesday showing that Americans ' opinions on former President Trump haven't changed during his ongoing hush money trial in New York City.
As closing arguments in the trial occurred, Enten cited a Quinnipiac poll showing that after the prosecution examined Michael Cohen, 46 \% of respondents believed Trump had done something illegal regarding payments to adult film star Stormy Daniels -- the same percentage as during pre-opening statements."The fact is that most Americans don't really care that much,"Enten told CNN anchor John Berman in a clip flagged by The Daily Caller, explaining why he believes public opinion hasn't changed.
CNN ' s Harry Enten discusses a Quinnipiac poll showing that New York ' s trial against Donald Trump has had no effect on Americans ' opinion on whether or not the former president is guilty of crimes @ @ @ @ @ @ @ @ @ @ in the"44 days"since jury selection began."Pre-opening statements, think Trump did something illegal, 46 \%, after the direct examination of Michael Cohen by the prosecution, look at where we are now, 46 \%,"Enten said."The percentage of Americans who think that the charges are very serious, in fact, dropped from 40 \% to a little bit more than 35 \% during the course of this trial."Former President Donald Trump, flanked by attorneys Todd Blanche and Emil Bove, arrives for his criminal trial at the Manhattan Criminal Court in New York, NY on Wednesday, May 29, 2024. ( Jabin Botsford/Pool via REUTERS ) The reporter then offered his opinion on why the needle hasn't shifted over the course of the widely publicized case."I think the question is, what ' s exactly cooking here?"he said."Why hasn't there been much of a change?
While folks like you and me, real news junkies, @ @ @ @ @ @ @ @ @ @ on, the fact is that most Americans don't really care that much."The reporter pulled up an Ipsos poll, showing that the public was more fixated on issues like the economy, immigration and abortion than Trump ' s time in court."The fact is John, when we ' re looking at these numbers, what we see is Americans ' minds aren't changing, and a big reason why Americans ' minds aren't changing is at this particular point, John, they are tuned out of the conversation,"Enten said.

% matched lemmas: conduit
\end{textsample}
